\chapter{The Rise of Stalin}

The death of Lenin brought into the open the issue which had long preoccupied the party leaders. Zinoviev had already assumed without hesitation the provisional mantle of the succession. Stalin had studiously refrained from disclosing his ambitions. At a commemorative session of the Union Congress of Soviets on January 26, 1924, the eve of the funeral, Stalin's tribute was distinguished from those of his colleagues by a fervent strain of worshipful devotion still unfamiliar in the Marxist or Bolshevik vocabulary: ``we communists'' were humble and loyal disciples, pledged to carry out every injunction of the dead master. Two notable decisions were taken. One was to re-name Petrograd ``Leningrad''; Lenin had superseded and eclipsed Peter in moulding the fortunes of the fatherland. The other was to strengthen the party by a mass recruitment of ``workers from the bench'', which was dubbed ``the Lenin enrolment''. The demand for a larger representation of workers in the party had figured in Trotsky's letter of October 8 and in the Politburo resolution of December 5, 1923 (see p. 66 above), and could be justified by much that Lenin himself had written. Its execution was now in the hands of Stalin, the general secretary of the party. 
% 注：去除了 "jwas"、".himself" 等边缘乱码，并修复了 "succes- sion" 和 "ambi- tions" 的跨行断字。

The Bolshevik party in 1917 had a membership of not more than 25,000. During the revolution and the civil war its numbers were progressively swelled by mass admissions. Statistics for this early period are unreliable. But early in 1921 it had reached a total of 600,000 or perhaps 700,000. The purge ordered by the tenth party congress in March 1921 was drastic. Some members, recruited in the enthusiasm of the revolution and the civil war, drifted away; others were expelled as unsuitable. By the beginning of 1924, the membership had been reduced to 350,000. The Lenin enrolment, which in two years added 240,000 new members to the party, increasing its numbers by more than two-thirds, was hailed both as a move towards more democracy and as an assertion of the rightful predominance of genuine workers in the party, though its later stages also included a substantial enrolment of peasants. Its historical role was quite different. It was the symbol of a gradual change in the character of the party which had deeper causes. Almost unnoticed a new conception emerged which differentiated the party of Stalin from the party of Lenin. 
% 注：修正了 "enrol- tnent" 为 "enrolment"，并去除了 "/the" 的斜杠。

Lenin before the revolution had envisaged the party as a small homogeneous group of devoted revolutionaries pledged to the overthrow of a regime of inequality and oppression. Even after the revolution, he continued to think of the party as an elite group of dedicated workers; and he was more concerned to purge the unfit than to open a wide door to recruitment. The sharp reduction in the number of party members between 1921 and 1924 was certainly due to his insistence. Lenin, though he had moved a long way from the Utopian views expressed in State and Revolution, still looked forward, in the words of the party programme of 1919, to ``a simplification of the functions of administration, accompanied by a rise in the cultural level of the workers'', and seemed unaware, till the very end of his life, of the vast complexities and problems of public administration. By this time the conception of an elite party was an anachronism. In 1920, 53 per cent of party members were said to be working in Soviet institutions of one kind or another, and 27 per cent were in the Red Army. Gradually and insensibly, the party had been transformed into a machine geared to conduct and supervise the affairs of a great state. It was the plain duty of rank-and-file members---and especially of new members who lacked the revolutionary grounding of the pre-1917 generation---to support the leaders loyally in this formidable task; and party membership carried with it certain undeclared privileges which made the performance of this duty worth while. The Lenin enrolment was accompanied by a further purge of undesirable members; and, since both the purge and the enrolment were controlled by the party secretariat, it may be guessed that adherence to the new party orthodoxy was one of the main criteria applied. The Lenin enrolment, and the whole process of which it formed part, enhanced the power of the party machine and of the general secretary who manipulated it. Molotov spoke no more than the truth when he observed at the party congress of 1924 that ``the development of the party in the future will undoubtedly be based on this Lenin enrolment''. 
% 注：去除了 "f‘a"、"[panied"、"Ithe"、"ithe"、"'certain"、"|).he"、"1 Informed"、"f//of" 等等大量的单侧标记符号。

Another and more subtle change followed the replacement of Lenin's elite party by the mass party of Stalin. The party statute imposed on members the obligation, once a policy decision had been taken, to speak in support of it with a united voice. Loyalty to the party meant acceptance of its discipline. But it was assumed that the decision would have been taken by democratic procedures after free discussion among party members. Nor did anyone suggest that the party was infallible; Lenin often drew attention to mistakes that had been made, and admitted errors of his own. When his fiftieth birthday was celebrated in April 1920, at the moment of victory in the civil war, he rather strangely spoke, in his reply to the greeting of his comrades, of the danger of the party ``giving itself airs''. The angry controversies which divided the party on the eve of the introduction of NEP shocked Lenin and other party leaders into a realization of the hazards involved in the unfettered toleration of dissent; and the Kronstadt mutiny increased the sense of alarm. The disciplinary measures adopted by the tenth congress were an ominous landmark in party history. But Lenin never made his peace with the conception of a central party organization announcing infallible edicts and imposing silence on all dissent within the party and outside it. When at the last party congress which he was to attend, in March 1922, he observed that the party had enough political and enough economic power, and that ``what is lacking is culture'', he already showed a troubled consciousness of the dangers that lay ahead. In the last tormented months of Lenin's active life he was preoccupied both by mistrust of Stalin's personality and by the need to struggle against ``bureaucracy'' in the party as well as in the state. Belief in the infallibility of the party, in the infallibility of Lenin, and eventually in the infallibility of Stalin himself, was a later development, the seeds of which were sown in the first weeks after Lenin's death. 
% 注：去除了 "•ill"、"///of"、"jjidecision"、"///a"、"/((its"、"IINEP"、"is peace" (修正为 his peace)、"—^H/announcing"、"^/('within" 等极端破碎的噪点字符。

While the Lenin enrolment was in progress, Stalin took a further step to distinguish himself as Lenin's most faithful disciple. He delivered at the Sverdlov University six lectures ``On the Foundations of Leninism'', which were published in Pravda. They were clear, schematic and entirely conventional. One sentence only might, in the light of subsequent developments, have attracted attention: 

\begin{quote}
For the final victory of socialism, for the organization of socialist production, the efforts of one country, particularly of a peasant country like Russia, are insufficient; for that, the efforts of the proletarians of several advanced countries are required. 
\end{quote}
% 注：使用 quote 环境来排版段落中的独立引言。

But this was merely the recital of a familiar item in the party creed. The lectures passed without comment. The other leaders showed no interest in Stalin's incursion into the field of theory where he had hitherto rarely sought to shine. What was significant in Stalin's initiative was the consecration of a specific cult of ``Leninism''. If the term was current during Lenin's lifetime, it was used, like ``Trotskyism'' later, as a term of opprobrium by opponents eager to discredit it. Henceforth, on the lips of Stalin and of other party leaders, Leninism was a vaguely defined, but infallible, body of doctrine, which distinguished the official party line from the heresies of its critics. 

The embarrassment of Lenin's testament had still to be overcome. Fortunately for Stalin, his embarrassment was shared by the other leaders, none of whom escaped unscathed. At what precise moment they became aware of its contents is not recorded. But on May 22, 1924, on the eve of the thirteenth party congress, a gathering of prominent party members heard it read by Kamenev, who presided. Then Zinoviev spoke, in terms of fulsome devotion to the dead leader, ending with the verdict that ``on one point'' Lenin's apprehensions had proved unfounded, and that it was not necessary to remove Stalin from his post. Kamenev supported Zinoviev. Nobody expressed any other view. Trotsky, just back from the Caucasus, sat silent through the proceedings. The only clash arose over Krupskaya's insistent demand that the testament should be read to the congress. The meeting decided, by a majority of 30 to 10, that it would be sufficient to communicate it confidentially to the leading delegates. 

The problem of the opposition loomed large at the congress. Zinoviev restrained himself in his main report, ending with a rhetorical appeal to members of the opposition to come to the tribune to confess their error, and admit that the party was right. Many delegates denounced the opposition, and Trotsky by name. Trotsky rose painfully and reluctantly to meet Zinoviev's challenge. ``One cannot be right against the party,'' he now proclaimed. The party could make ``particular mistakes''; and he continued to believe that the resolution of the January conference condemning him was ``incorrect and unjust''. Nevertheless as a loyal party member he was bound to say: ``Just or unjust, this is my party, and I bear the consequences of its decision to the end.'' Whether one regards it as the source of the inhibition which prevented Trotsky from giving battle, or as the rationalization of an inhibition which had deeper psychological roots, this declaration of submission, coupled with a refusal to confess error, was significant for his attitude at this time. It was only two years later---when it was already too late---that he regained his freedom of action, struck boldly at his enemies, and rallied his friends for his defence. The congress heard a plea from Krupskaya for peace between the factions and for ``an end to further discussion''. This went unheeded. Stalin and Zinoviev wound up the proceedings with speeches full of vituperation of Trotsky. He was, however, re-elected to the party central committee---apparently by a narrow margin. It is said that Zinoviev and Kamenev sought to exclude Trotsky from the Politburo, but that the proposal foundered on the opposition of Stalin, anxious to preserve his reputation as a moderate. 
% 注：修复了多行标点错位 "[Kvas", "Itwo", "/his"，重组了语序。

During the rest of the year Trotsky's literary prowess added fuel to the flames. In a commemorative pamphlet On Lenin he described his close personal association with Lenin at the time of the revolution in terms which seemed to inflate his own importance, and to relegate other participants to a secondary place. In October 1924, he published a long essay entitled Lessons of October, pillorying Kamenev and other ``old Bolsheviks'' for their resistance to Lenin's April theses on Lenin's return to Petrograd in April 1917, and the opposition of Zinoviev and Kamenev to the seizure of power in October, which had been cited by Lenin in the testament, with the qualification equating it with Trotsky's non-Bolshevik record as things which should not be brought up against them (see p. 62 above). This onslaught provoked a spate of controversial replies, and encouraged the triumvirate and their followers to delve, deeply and spitefully, into Trotsky's own record. Kamenev delivered a lengthy speech, published as a pamphlet under the title Leninism or Trotskyism?, in which he accused Trotsky of Menshevism, recounted his many sharp exchanges with Lenin, and added the henceforth familiar charge of ``under-estimating the peasantry''. Stalin followed, more briefly and incisively, in the same vein. Denunciation of Trotsky became a routine exercise in the press and in party meetings. The most savage blow was the discovery and publication of a forgotten letter written by Trotsky in 1913, full of crude and angry invective against Lenin. No further evidence was required to prove the incompatibility of ``Trotskyism'' with ``Leninism''. 
% 注：在此段的后半部分，OCR 扫描将 “Denunciation of Trotsky... 和 I .and publication of...” 严重交叉在了一起，并带有 "11913" 的数字错误。现已完全依照英语逻辑将其复位。

Overwhelmed by this flood of invective Trotsky remained silent. He succumbed once more to the mysterious malady which had afflicted him in the previous winter, and doctors advised his removal to a milder climate. He did not attend the session of the party central committee in January 1925. He addressed a letter to it in which he claimed that his silence in face of ``many untrue and even monstrous charges'' had been ``correct from the standpoint of the general interests of the party''; and ``in the interest of our cause'' he asked to be released from his duties as president of the Military-Revolutionary Council. He left for the Caucasus while the session was in progress. The committee hesitated over what sanctions to apply to him. Extremists, who included Zinoviev and the Leningrad delegation, proposed to expel him from the party, from the party central committee, or at the very least from the Politburo. Moderates, supported by Stalin, were content to relieve him of his military functions. The latter view prevailed; Trotsky was removed from the post of president of the Military-Revolutionary Council and People's Commissar for War. He was succeeded by Frunze, whose appointment was the signal for a powerful campaign to build up the Red Army, neglected since the end of the civil war. 

The controversy provoked by Lessons of October led, almost casually, to an important innovation in party doctrine. One of the items on which Lenin and Trotsky had once differed, and which was now brought up against Trotsky by his critics, was the so-called theory of ``permanent revolution''---a phrase originally used by Marx. Trotsky in 1905 argued that a revolution breaking out in backward Russia, while in its first stage it would remain a bourgeois anti-feudal revolution, would automatically pass over into the stage of a socialist anti-capitalist revolution. Lenin was unwilling to accept the prospect of this transition unless, as both he and Trotsky expected, revolution in Russia kindled the flame of revolution in the advanced countries of the west. The dispute was of little importance, and had been forgotten long before 1917, when Lenin in his April theses appeared to adopt a position nearer to that of Trotsky. But nobody displayed any interest in the issue till Bukharin, in December 1924, made his contribution to the campaign against Trotsky in an article on ``The Theory of Permanent Revolution''. Bukharin was concerned merely to spotlight Lenin's dissent from Trotsky, and drew no positive conclusion. But when, a few days later, Stalin also published a lengthy essay on the theme, written as an introduction to a collection of his speeches and articles, he used his denunciation of Trotsky's theory as the springboard for a new doctrine of ``socialism in one country''. 
% 注：去除了 "I " 的乱码标记。

Stalin now abandoned what he later called ``the incomplete and therefore incorrect'' formula in his lectures of the previous spring, in which he had held that the efforts of one country were ``insufficient for the organization of socialism''. Having declared that ``Trotsky's `permanent revolution' is the negation of Lenin's theory of proletarian revolution'', he proceeded to argue that Lenin had in several passages in his writings contemplated the possibility of a victory of socialism in one country. Stalin admitted that ``for a complete victory of socialism, for a complete guarantee against a restoration of the old order of things, the combined efforts of the proletariat of several countries are indispensable''. But did this mean that ``revolutionary Russia could not stand up against conservative Europe'', and build a socialist regime in the USSR? 
% 注：修正了 "Stahn" 这一明显的字母异化错误为 "Stalin"。

Stalin's answer was a resounding negative. The argument was complicated and casuistical, resting extensively on quotations taken out of context. It was also somewhat unreal, since it was conducted in conditions which neither Lenin nor Trotsky had considered possible---the survival of the revolutionary regime in Russia in the absence of revolution in other countries. But psychologically its impact was enormous. It supplied a positive and definable goal. It dispensed with vain expectations of help from abroad. It flattered national pride by presenting the revolution as a specifically Russian achievement, and the building of socialism as a lofty task in the fulfilment of which the Russian proletariat would set an example to the world. Hitherto the dependence of the prospect for socialism in Russia on socialist revolution in other countries had occupied a central place in party doctrine. Now the order of priority was reversed. Stalin boasted that the victory of the revolution in Russia was ``the beginning and the premiss of world revolution''. Critics of Stalin's doctrine were, implicitly and explicitly, revealed as faint-hearted, timid and mistrustful of the Russian people, sceptical of their capacity and determination. Socialism in one country was a powerful appeal to national patriotism. Indisputably it put Russia first. 
% 注：去除了 "ftrine" 和 "/a" 的杂乱符号。

Stalin had created a climate of opinion which he was to exploit to the utmost in his struggle against his rivals. But for the moment nobody took his abstruse excursion into theory very seriously. At the session of the party central committee in January 1925, which condemned Trotsky, socialism in one country was not mentioned. Bukharin hesitatingly reverted to it in a speech three months later, without mentioning Stalin, and in terms which suggested that he himself was one of its authors. It appeared, not very prominently, in the main resolution of the party conference in April 1925, which, on the strength of quotations from Lenin, announced that ``in general the victory of socialism (not in the sense of final victory) is unconditionally possible in one country''. When the triumvirate broke up some months later, it was alleged that this passage had been the subject of a clash in the Politburo on the eve of the conference. But the evidence suggests that Zinoviev and Kamenev raised no strong objection, and were indifferent rather than hostile. When Stalin celebrated this modest victory in a speech after the conference, he produced yet another quotation from Lenin: 

\begin{quote}
Only when the country is electrified, only when industry, agriculture and transport have been put on the technical basis of modern large-scale industry, only then shall we be finally victorious. 
\end{quote}

Hitherto socialism in one country might have been seen as a continuation of NEP, which had also turned its back on the bleak prospect of international revolution, and marked out the road to socialism through an alliance with the Russian peasant. Now Stalin was groping his way to the very different conception of a self-sufficient Russia, transformed and rendered economically independent through a modernized industry and agriculture. Stalin did not press the point, and was perhaps not yet fully conscious of its implications. But it was a dazzling long-term vision; and it fitted in with changes which were beginning to make themselves felt on the economic scene. 
% 注：本段修复了最严重的字符异化和分栏倒置。修复包括 "I me bleak prospect", "internafional feVoiufion" (修复为 international revolution), "marked put" (修复为 marked out), "\yith", "/ypeasant", "I llconception" 等。

The gradual rise of Stalin to a position of authority after the death of Lenin occurred in a period of acute economic controversy and conflict, which was also a period of economic revival. The resolution of the scissors crisis in December 1923 and subsequent party pronouncements heralded a new attention to the restoration of heavy industry. The doctrine of socialism in one country, whatever the intention of its exponents, lent support to the promotion of heavy industry as a condition of self-sufficiency. But it also implied that this could be achieved with the resources of the backward Russian economy. Here lay the crux. The controversy about industrialization, like every other issue in the Soviet economic scene, was bound up with the problems of agriculture, which once more disturbed the current mood of complacency. The grain harvest of 1924, though impaired by a late summer drought, was fair. Nobody seems to have doubted that the peasant, freed from the burdens of the scissors crisis, would deliver to the state collecting organs, at officially fixed prices, the quantities of grain required to feed the cities. Nothing of the kind happened. The grain collections fell disastrously short. 
% 注：这是另一个被 OCR 严重切割破坏的超长段落 ("...»||the death... Ijtion... /|controversy... y: kocialism... inteiTtion... (peasant... ")，我已通过重新拼接，完全恢复了其连续性和逻辑性。

Private traders appeared for the first time in large numbers on the market, and the fixed prices had to be abandoned. At the turn of the year prices were rising fast. The price of rye doubled between December 1924 and May 1925. With the return of the free market, the scissors had once more opened, this time in favour of the peasant, and the cities were held to ransom. Moreover, the price mechanism operated to increase disparities of wealth in the countryside. It was the rich peasant, the hated kulak, who had large surpluses to sell, and could afford to hold them till prices reached their peak. It was reported that many poor peasants, pressed to realize their crop, had sold cheap in the autumn to kulaks, who made their profits by selling dear in the spring. 

These developments were the starting-point of a keen controversy in the party. The leaders clung to the guiding principle of NEP---conciliation of the peasant; Zinoviev in July 1924 had launched the slogan ``Face to the Countryside''. A few days later, Preobrazhensky read to the Communist Academy a paper on ``The Fundamental Law of Socialist Accumulation'', which was widely recognized as a searching challenge to the official line. Marx had shown that the early stages of capitalist accumulation had required ``the separation of the producers from the means of production'', i.e. the expropriation of the peasantry; so---argued Preobrazhensky---socialist accumulation ``cannot do without the exploitation of small-scale production, without the expropriation of part of the surplus product of the countryside and of artisan labor''. He rejected as impracticable the principle of ``equivalent exchange'' between town and country, and advocated ``a price policy consciously directed to the exploitation of the private economy in all its forms''. Preobrazhensky did not mince his words; and his outspokenness gave a handle to defenders of the party leadership and of the peasantry. Bukharin published an indignant reply, which denounced the article as ``the economic foundation of Trotskyism''. But Preobrazhensky had confronted the party in the plainest terms with the harsh dilemma of reconciling the process of industrialization with continued indulgence for the peasant. 
% 注：去除了段首错落的 "[The", "Ireached", "ipriation", "[accumulation" 并修复了段落结构。

Throughout 1925, while Stalin manoeuvred cunningly between the other leaders, an open clash between the two policies was avoided. Pressure was strong for further concessions to the peasant, which meant in practice concessions to the well-to-do peasant or kulak. A party conference in April 1925 voted three such measures. The agricultural tax, the sole vehicle of direct taxation in the countryside, was to be reduced, and its incidence changed to make it less progressive. The right to employ hired labour, and the right to acquire land by leasing, still partially though ineffectively restricted by the agrarian code, were to be recognized. It was at this moment that Bukharin delivered a speech which was long quoted as the most outspoken exposition of the policy represented in these decisions. He pleaded the cause of ``the well-to-do top layer of the peasantry---the kulak and in part the middle peasant'', who needed incentives to produce. ``To the peasants, to all the peasants'', he exclaimed, ``we must say: Enrich yourselves, develop your farms, and do not fear that constraint will be put on you''. He denied that this was ``a wager on the kulak'' (a phrase coined fifteen years earlier to describe the Stolypin reform). But he equally rejected ``a sharpening of the class war in the countryside''. Bukharin, like his adversary Preobrazhensky, compromised his cause by undiplomatic frankness. Stalin seems to have told other party leaders that ``Enrich yourselves'' was ``not our slogan''. But it was some months before it was publicly disowned; and the course charted by Bukharin was followed during the rest of the year. 
% 注：去除了 "/? It^"、"Ijto"、"(well-to-do"、"jvehicle" 等极多的边缘乱码。修复了 ";o-do" 为 "-to-do"。

Side by side, however, with policies designed to provide incentives for peasant production, the needs of heavy industry attracted increasing attention. Hitherto the revival of industry had for the most part meant the bringing back into productive use of plant and machinery which had been idle since the civil war; for this no great capital outlays were required. But by the end of 1924 the process had reached its limit. It was estimated that existing factories and plant were being utilized to 85 per cent of capacity. Industry was beginning to approach the levels of production achieved in 1913, and could contemplate an advance beyond them. But to maintain the rate of industrial growth, and especially to revive heavy industry, demanded large-scale capital investment. The party central committee in January 1925 advocated ``budget allocations'' to industry as well as ``an expansion of credit''. Obsolescent equipment must be renewed, and new industries created. Thus encouraged, Vesenkha convened ``a special conference on the restoration of fixed capital in industry'', which remained active for the next 18 months. The party conference of April 1925, which voted concessions to the peasant, also approved a three-year plan for the metal industry, involving a total investment of 350 million rubles. 

The year 1925 was still a period of optimism, when it seemed possible to meet all demands of a rising economy. It was not the harvest itself, the best since the revolution, but the sequel of the harvest, which in the last months of the year showed up the dimensions of the problem inherent in relations between industry and agriculture. The state grain-collecting organs abandoned the ``fixed'' prices of 1924, and were instructed to work on ``directive'' prices which could be adjusted from time to time. In spite of the experience of the previous harvest, everyone seems to have assumed that the abundance of the crop would keep prices down, that surplus grain would be available for export, and that the proceeds of the harvest would provide funds for the financing of industry. These hopes were disappointed. After the harvest of 1925 the prosperous peasants held large stocks of grain. But they had no incentive to exchange these for money. The reduction of the agricultural tax had relieved the pressure of taxation; the supply of industrial goods was meagre, and offered little that they wished to buy; and, though the currency had nominally been stabilized, a hoard of grain was a more reliable asset than a packet of bank notes. They could afford to wait. Grain came slowly to the market. Under the influences of short supplies, competition with free-market purchasers, and even competition between different state purchasing organs, prices soared. Hopes of grain exports, or of profits from the harvest to finance industry, evaporated. The harvest had been a success for the peasant. The marketing of the harvest was a disaster for the government. The crisis split the party, and gave the signal for a prolonged and bitter struggle between the claims of industrialization and planning on the one hand, and the peasant-oriented market economy promoted by NEP on the other, which was to dominate the ensuing period. 

These events were the background of Stalin's ascent to a position of supreme authority in the party and in the USSR. The year 1925 was decisive. Fear and jealousy of Trotsky was the cement which held the triumvirate together. After his defeat and demotion in January 1925 it began slowly to crumble. Trotsky spent more than three months convalescing in the south. When he returned to Moscow at the end of April 1925 he faced an embarrassing situation. Eastman, a well-known American communist, had spent the winter of 1923-1924 in Moscow, and was a declared supporter of Trotsky. At the beginning of 1925 he published in New York a small book Since Lenin Died which gave a detailed and accurate account, from Trotsky's standpoint, of the intrigues conducted by the triumvirate during the last weeks of Lenin's life and after his death, and quoted Lenin's testament---the first reference to that document to appear anywhere in print. The revelations made a sensation. Anxious members of the British Communist Party wrote and telegraphed to Trotsky asking for his comments for publication. The party leaders in Moscow insistently demanded from him a refutation of Eastman's indictment. Trotsky was once more faced with the dilemma of standing his ground, or refusing to fight on what might be called a secondary issue. He still suffered from the deep inhibition which prevented him from coming out publicly in opposition to a majority of his colleagues: ``one cannot be right against the party''. If it occurred to him that to retreat meant to compromise his cause and disavow his friends, he stifled these doubts in the name of party discipline. On July 1, 1925, he signed a long statement which, as he wrote three years later, ``was forced on me by a majority of the Politburo''. He described the allegation that the party central committee had ```concealed' from the party a number of extremely important documents written by Lenin in the last period of his life'', including ``the so-called `testament''', as ``a slander''. Lenin had left no testament; everything that he had written, notably ``one of Vladimir Ilich's letters containing advice of an organizational character'', had been communicated to the delegates at the party congress. Stories of a concealed testament were ``a malicious invention''. Trotsky's statement was published in the British Left-wing Sunday Worker on July 19, and in the Russian party journal Bolshevik on September 1, 1925. It was the last triumph of the united triumvirate. 
% 注：去除了段首杂乱的逗号和单侧括号。

Trotsky on his return to Moscow had been appointed to two or three minor, and largely nominal, posts connected with industry. During the rest of the year he made a few speeches, and wrote a few articles, on industrial development and planning, stressing the need to ``catch up with the west'', but offering no direct challenge to party policy. His restraint relaxed the last bond that held the triumvirate together. After some preliminary bickering, open dissension broke out over the grain collections crisis. Zinoviev and Kamenev, reversing their earlier position, came out against the peasant orientation, of which Bukharin remained the most articulate champion. In September Zinoviev submitted for publication in Pravda an article entitled ``The Philosophy of an Epoch''. It took the form of an attack on an emigre writer, Ustryalov, who had enthusiastically endorsed Bukharin's support of the kulak, and joyfully proclaimed that ``the peasant is becoming the only real master of the Soviet land''. Zinoviev concluded that ``NEP, together with the delay in world revolution, is really pregnant, among other dangers, with the danger of degeneration''. The party central committee insisted on the removal of phrases that pointed too directly at Bukharin. But the sense of the article, which ran through two issues of Pravda, and was published as a pamphlet, could not be mistaken. In the following month Zinoviev published a volume of essays under the title Leninism. One of these repeated the attack on Ustryalov, and denounced the slogan ``Enrich yourselves'', though still without mentioning Bukharin by name. Another quoted Lenin's denunciation of the kulaks, and recalled his description of NEP as a ``retreat''; this implied that Soviet industry under NEP was a form of ``state capitalism''---a conclusion denied by Bukharin. The most decisive chapter of all was a frontal attack on ``socialism in one country''; it was impossible to ``remain Leninists if we weaken by a single jot the international factor in Leninism''. This was a declaration of war, not only against Bukharin, but against Stalin himself. 

Zinoviev's abrupt abandonment of the peasant orientation, and espousal of the cause of industry and the proletariat, had a certain logic. A struggle for mastery between Zinoviev and Stalin was a struggle between the Leningrad party organization controlled by the former and the central party organization in Moscow dominated by the latter. Kamenev was head of the local Moscow organization. But this was overshadowed by the central organization in the same city; Kamenev lacked the authority to assert its independence, and was soon ousted. Leningrad was still the most heavily industrialized city in the USSR. It was the home of the proletariat which had been the vanguard of the revolution, and maintained its proletarian tradition. In Moscow the new proletariat had retained far closer ties with the countryside. Zinoviev could mobilize and lead the Leningrad workers against Moscow only from a platform which upheld the pre-eminent claims of the workers, and rejected with contumely the attempt to exalt the role of the peasant. Rivalry between the two capitals, and between the two party organizations, between Pravda, which was the organ of the party central committee in Moscow, and Leningradskaya Pravda, the newspaper of the Leningrad party organization, played a significant part in the struggle for power between Stalin and Zinoviev. 

The battle-ground was the fourteenth party congress which sat during the last fortnight of 1925. Stalin and Zinoviev were the main speakers; Bukharin answered Zinoviev, and was answered by Kamenev. While Zinoviev and Kamenev fiercely denounced the kulak, Bukharin stood his ground; and Stalin, whose concern was to defeat his two principal rivals, half-heartedly supported Bukharin. The congress took no significant decision on agricultural policy. But it registered a growing impatience with the favours enjoyed by the kulak, and it harped once more on the urgency of industrialization. When the dust of party conflict had settled, it was clear that a major decision lay ahead. Bukharin at the congress, in a phrase which was long remembered, made a desperate attempt to prove that conciliation of the peasant was not incompatible with policies of industrialization: ``We shall move forward at a snail's pace, but all the same we shall be building socialism, and we shall build it''. But snail's pace industrialization no longer satisfied the growing body of opinion which wanted to transform the USSR into a great industrial country, independent of the west. Paradoxically the victory of Bukharin and the defeat of Zinoviev at the congress did not lead to the victory and defeat of the policies for which they respectively stood. It was not altogether misleading when the congress was later dubbed ``the congress of industrialization''. 

Economic problems did not, however, dominate the debate, which opened in a fairly low key, but became more acrimonious when sensitive political and personal issues were broached. Kamenev criticized ``the theory of a `leader''', and launched a personal attack on Stalin. Krupskaya spoke for the opposition, and made a sensation by challenging the doctrine that ``the majority is always right''. Molotov and Mikoyan were among those who supported the official line, and Voroshilov eulogized Stalin. The delegates on both sides, nominally elected by their party constituencies, had been hand-picked by the party organizations, and a solid phalanx of Leningraders was isolated in a hostile audience. The resolution endorsing the official line was voted by a majority of 559 to 65. Leningradskaya Pravda, hitherto the mouthpiece of the opposition, was taken over, a new editor being appointed from Moscow. After the congress a powerful delegation, which included Molotov, Voroshilov, Kalinin, Rykov, Tomsky, Kirov, and later Bukharin, proceeded to Leningrad, and addressed a series of mass meetings of party members. The means of pressure which had silenced and intimidated the followers of Trotsky were now brought to bear on the supporters of Zinoviev. Mass meetings of workers were induced, by large majorities, to condemn their old leaders and vote approval of the congress decisions. The ground was thus prepared for a Leningrad provincial party conference, at which Bukharin was the principal speaker. The same verdict was repeated, and loyal supporters of the party central committee were elected to the Leningrad party organs; Kirov, a young and popular recruit to the party leadership, became secretary of the Leningrad provincial party committee, the de facto head of the Leningrad organization. It was a complete take-over. Zinoviev remained a member of the Politburo and president of Comintern. But, expelled from his base in Leningrad, he lost all effective power. Stalin was the victor. But what his victory portended, either economically or politically, was still unclear.
% 注：去除了段末出现的 "ius" 字符。